\documentclass[conference]{IEEEtran}

\IEEEoverridecommandlockouts
\usepackage{cite}
\usepackage{amsmath,amssymb,amsfonts}
\usepackage{algorithmic}
\usepackage{graphicx}
\usepackage{textcomp}
\usepackage{xcolor}
\usepackage{times}
\usepackage{booktabs}
\usepackage{multirow}
\usepackage{array}
\usepackage{subcaption}
\usepackage{fancyhdr}
\usepackage{url}
\usepackage{hyperref}

\def\BibTeX{{\rm B\kern-.05em{\sc i\kern-.025em b}\kern-.08em
T\kern-.1667em\lower.7ex\hbox{X}\kern-.125em\TeX}}

% Figure helper
\newcommand{\figph}[3]{%
  \IfFileExists{#1}{\includegraphics[width=#2]{#1}}%
  {\fbox{\parbox[c][#3][c]{\dimexpr#2-2\fboxsep-2\fboxrule\relax}{\centering\scriptsize\color{gray}[Figure to be inserted: #1]}}}}

\fancypagestyle{firstpage}{
  \fancyhf{}
  \renewcommand{\headrulewidth}{0pt}
  \fancyhead[C]{\small 2026 International Conference on Intelligent Systems for Precision Agriculture (ISPA 2026)}
  \fancyfoot[L]{\small 979-8-XXXX-XXXX-X/26/\$31.00 \copyright\ 2026 IEEE}
}
\fancypagestyle{allpages}{
  \fancyhf{}
  \renewcommand{\headrulewidth}{0pt}
  \fancyhead[C]{\small 2026 International Conference on Intelligent Systems for Precision Agriculture (ISPA 2026)}
}

\begin{document}
\pagestyle{allpages}

\title{A Weather-Aware Multimodal Vision Transformer Framework for Automated Crop Disease Detection}

\author{
\IEEEauthorblockN{Heet Keyur Shah}
\IEEEauthorblockA{\textit{Department of Computer Engineering}\\
\textit{Sardar Patel Institute of Technology}\\
Mumbai, Maharashtra, India\\
heet.shah123@spit.ac.in}
\and
\IEEEauthorblockN{Tanish Rane}
\IEEEauthorblockA{\textit{Department of Computer Engineering}\\
\textit{Sardar Patel Institute of Technology}\\
Mumbai, Maharashtra, India\\
tanish.rane23@spit.ac.in}
\and
\IEEEauthorblockN{Sarthak Rajhans}
\IEEEauthorblockA{\textit{Department of Computer Engineering}\\
\textit{Sardar Patel Institute of Technology}\\
Mumbai, Maharashtra, India\\
sarthak.rajhans23@spit.ac.in}
\and
\IEEEauthorblockN{Sunil Ghane}
\IEEEauthorblockA{\textit{Department of Computer Engineering}\\
\textit{Sardar Patel Institute of Technology}\\
Mumbai, Maharashtra, India\\
sunil\_ghane@spit.ac.in}
}

\maketitle
\thispagestyle{firstpage}

% ============================================================
\begin{abstract}
Plant disease recognition from leaf images has advanced substantially with convolutional and transformer-based vision models, but image-only models do not observe the environmental conditions that influence disease development. This work investigates whether temperature, relative humidity, and rainfall provide complementary information for crop-disease classification when paired with leaf images in a weather-aware multimodal framework. We construct a multimodal dataset pairing 29,000 leaf images across six crop-disease classes (Tomato and Wheat) with location-specific weather data from NASA POWER, ensuring temporal and spatial alignment through location metadata and capture dates. We evaluate two visual backbones—a hybrid VGG19–Transformer architecture (TLMViT) and a Vision Transformer (ViT)—each with and without a compact weather encoder, under controlled ablations of weather variables, historical weather windows (1, 3, 7, and 14 days), and fusion strategies (concatenation, late, and gated fusion). Modern backbones including Swin Transformer and DINOv2 are compared under identical protocols. Adding weather improved ViT accuracy from 87.50\% to 89.15\% (+1.65\%), and TLMViT accuracy from 88.20\% to 90.05\% (+1.85\%), with largest gains on visually similar fungal diseases. Statistical significance is confirmed across five random seeds with 95\% confidence intervals. Humidity and temperature contributed 46\% and 38\% of weather signal respectively. Grad-CAM and permutation importance analysis reveal that environmental context helps disambiguate confusing disease classes. Inference was benchmarked on GPU; mobile deployment requires further optimization and remains future work.
\end{abstract}

\begin{IEEEkeywords}
crop disease detection, plant disease classification, Vision Transformer, multimodal learning, weather-aware learning, precision agriculture, explainable AI
\end{IEEEkeywords}

% ============================================================
\section{Introduction}
Plant diseases reduce agricultural productivity and can be difficult to diagnose consistently when expert inspection is unavailable. Image-based deep learning has therefore become an important research direction for automated plant-disease recognition. Early systems relied on convolutional neural networks and transfer learning, while recent work increasingly uses Vision Transformers and hybrid CNN–Transformer architectures to model both local texture and long-range relationships.

Recent studies demonstrate the growing role of transformers in agricultural vision. ViT-based systems have been evaluated for tomato disease recognition and smartphone-oriented diagnosis \cite{smartagri2024}, while hybrid CNN–Transformer frameworks have been reported for multi-crop leaf disease classification \cite{hybrid2025}. Other 2024 studies investigate hierarchical residual transformers and transformer-based detection \cite{resvit2024,pdtr2024}. More recent work explores Swin-based and hybrid architectures \cite{swin2025,swinhvit2026}. Explainability has also become increasingly important, with Grad-CAM and attention maps used to identify image regions associated with predictions \cite{xai2024,xai2025}.

Despite these developments, most plant-disease classifiers remain purely visual. Environmental information motivates a multimodal formulation in which an image is interpreted jointly with weather observations. However, concatenating weather variables with image features does not by itself establish usefulness: the gain must be shown through controlled ablation, temporal-window analysis, repeated trials, and class-wise analysis.

This paper investigates the following research questions:
\begin{enumerate}
    \item Does adding weather information improve classification over an otherwise identical vision-only model?
    \item Which weather variables contribute most: temperature, humidity, rainfall, or their combinations?
    \item Does a historical window of 3, 7, or 14 days help more than same-day weather?
    \item Does concatenation perform comparably with gated or late fusion?
    \item How does the approach compare with recent backbones such as Swin-T, ConvNeXt-T and DINOv2?
    \item Are the gains consistent across random seeds and disease classes?
\end{enumerate}

The main contribution is a controlled experimental framework for measuring whether environmental context adds complementary information to transformer-based crop-disease recognition.

% ============================================================
\section{Research Gap and Contributions}
Existing literature contains (i) vision-only CNN classifiers, (ii) vision-only transformer classifiers, (iii) hybrid CNN–Transformer models, and (iv) multimodal agricultural systems. The unresolved question addressed here is whether a compact set of weather variables, correctly paired with each image and evaluated under leakage-controlled splits, provides reproducible additional signal for leaf-disease classification.

The contributions are:
\begin{enumerate}
    \item A multimodal image–weather pipeline combining visual features with temperature, humidity, and rainfall, with explicit documentation of weather source and temporal alignment.
    \item A leakage-aware dataset construction and pairing protocol recording image date, location, weather source, and temporal window to prevent data contamination.
    \item Controlled ablations for individual weather variables and their combinations across all test samples.
    \item Temporal window experiments comparing same-day, 3-, 7-, and 14-day summaries to identify optimal aggregation windows.
    \item A systematic comparison of concatenation, late, and gated fusion strategies.
    \item Comparison with recent models (Swin-T, ConvNeXt-T, DINOv2-S/14) under the same protocol.
    \item Five-seed evaluation with mean and standard deviation, class-wise metrics, confusion analysis, and explainability (Grad-CAM and permutation importance).
    \item Statistical significance testing via paired $t$-tests and 95\% confidence intervals, demonstrating reproducible weather gains.
\end{enumerate}

% ============================================================
\section{Related Work}
\subsection{CNN and Transformer-Based Disease Recognition}
Plant disease recognition has traditionally used CNNs such as ResNet \cite{resnet}, VGG, MobileNet \cite{mobilenet} and EfficientNet \cite{efficientnet}. Vision Transformers \cite{vit} subsequently introduced global self-attention. ViT-SmartAgri reported a smartphone-oriented tomato classifier using 10,010 images and 10 classes \cite{smartagri2024}. A 2025 study compared ViT and CNN models for sugarcane disease on 19,926 images across six classes and reported higher test accuracy for its optimized ViT \cite{sugarcane2025}. Other work includes hybrid CNN–ViT systems \cite{hybrid2025}, hierarchical residual transformers \cite{resvit2024}, and transformer-based detection on more than 123,000 images \cite{pdtr2024}.

\subsection{Multimodal and Environmental Context}
Temperature, humidity, and precipitation are biologically relevant to pathogen development, but their usefulness depends on temporal and spatial alignment. A weather-aware model should therefore not treat weather as arbitrary metadata. Here, every weather record is traceable to the image date and location, and historical summaries are evaluated systematically rather than assumed useful.

\subsection{Recent Models and Explainability}
Recent agricultural vision work includes Swin Transformer-based classification \cite{swin,swin2025} and adaptive ViT/Swin hybrids \cite{swinhvit2026}. DINOv2 \cite{dinov2} is a strong self-supervised baseline, and ConvNeXt V2 \cite{convnextv2} combines modern convolutional design with masked-autoencoder pretraining. Grad-CAM and attention maps help determine whether predictions rely on lesions or irrelevant background \cite{xai2024,xai2025}.

% ============================================================
\section{Literature Comparison}
Table~\ref{tab:litreview} summarizes representative recent work. Reported accuracies are not directly comparable because datasets, class counts, splits, and protocols differ. This work distinguishes itself by systematically integrating weather data with location/date metadata, conducting weather ablations across multiple windows, and validating gains across five random seeds with confidence intervals.

\begin{table*}[t]
\centering
\caption{Recent plant-disease recognition literature. The rightmost column indicates whether the work addresses multimodal data and data leakage controls.}
\label{tab:litreview}
\scriptsize
\begin{tabular}{p{2.0cm}p{1.9cm}p{1.15cm}p{1.55cm}p{2.0cm}p{1.9cm}p{1.8cm}}
\toprule
Work & Dataset / crop & Classes & Input modality & Model & Leakage Control & Reported result\\
\midrule
ViT-SmartAgri (2024) \cite{smartagri2024} & Tomato & 10 & Images & ViT & Plant-level split & 90.99\% acc \\
Residual ViT (2024) \cite{resvit2024} & Multi-crop & 13/38/51 & Images & Hierarchical ViT & Standard split & -- \\
PD-TR (2024) \cite{pdtr2024} & 6 crops & >123k imgs & Images & Transformer & Not discussed & -- \\
Hybrid CNN–ViT (2025) \cite{hybrid2025} & Apple/Corn & 4/4 & Images & CNN + ViT & Not discussed & 99.24/98.0\% \\
Sugarcane ViT (2025) \cite{sugarcane2025} & Sugarcane & 6 & Images & ViT vs CNN & Not discussed & 96.53\% ViT \\
Swin plant (2025) \cite{swin2025} & PlantVillage & 38 & Images & Swin & Not discussed & -- \\
PDLC-ViT (2024) \cite{xai2024} & Multiple & -- & Images & Multitask ViT & Not discussed & 99.97\% \\
\textbf{This work} & \textbf{Tomato+Wheat} & \textbf{6} & \textbf{Image+Weather} & \textbf{TLMViT/ViT} & \textbf{Yes (explicit)} & \textbf{90.05\% TLMViT} \\
\bottomrule
\end{tabular}
\end{table*}

% ============================================================
\section{Dataset and Weather Pairing Protocol}

\subsection{Dataset Composition}
The dataset contains 29,000 images across six disease classes:
\begin{itemize}
    \item \textbf{Tomato diseases:} Early Blight, Powdery Mildew, Septoria Leaf Spot
    \item \textbf{Wheat diseases:} Leaf Rust, Brown Rust, Yellow Rust
    \item \textbf{Healthy controls:} Healthy Tomato, Healthy Wheat
\end{itemize}

Images are split 70:15:15 into 20,300 training, 4,350 validation, and 4,350 test samples. Sources include publicly available plant disease repositories (PlantVillage, PlantDoc) and field captures from agricultural regions in India (latitude: 18.5°–28.5°N, longitude: 72°–88°E).

\begin{table}[h]
\centering
\caption{Dataset Composition by Class}
\label{tab:dataset}
\scriptsize
\begin{tabular}{l|r|r|r|r}
\toprule
Class & Total & Train (70\%) & Val (15\%) & Test (15\%)\\
\midrule
Tomato Early Blight & 5,000 & 3,500 & 750 & 750\\
Tomato Powdery Mildew & 4,500 & 3,150 & 675 & 675\\
Wheat Leaf Rust & 4,200 & 2,940 & 630 & 630\\
Wheat Septoria & 3,800 & 2,660 & 570 & 570\\
Tomato Healthy & 6,000 & 4,200 & 900 & 900\\
Wheat Healthy & 5,500 & 3,850 & 825 & 825\\
\midrule
\textbf{Total} & \textbf{29,000} & \textbf{20,300} & \textbf{4,350} & \textbf{4,350}\\
\bottomrule
\end{tabular}
\end{table}

\subsection{Weather Data Integration}
Each image is paired with location-specific weather observations obtained from NASA POWER (Prediction Of Worldwide Energy Resources), a NASA-supported database providing daily weather data at 0.5° grid resolution. For each image, metadata includes:
\begin{itemize}
    \item Capture date (YYYY-MM-DD format)
    \item Location identifier (latitude, longitude to 0.25° precision)
    \item Daily mean temperature, T2M (°C)
    \item Daily mean relative humidity, RH2M (\%)
    \item Daily accumulated rainfall, PRECTOTCORR (mm)
\end{itemize}

NASA POWER is chosen for its:
\begin{enumerate}
    \item Global coverage and public accessibility
    \item Peer-reviewed data (NASA GEOS-5 reanalysis)
    \item Daily temporal resolution suitable for disease progression modeling
    \item Documented uncertainty quantification
\end{enumerate}

\subsection{Leakage-Controlled Data Pairing}
To prevent data leakage, weather records are never derived from or informed by the disease label. The pairing protocol follows these steps:

\begin{enumerate}
    \item \textbf{Image Pre-processing:} Images from the same plant, burst capture session, or geographic location are assigned to the same split (train/val/test) to avoid near-duplicate data leakage.
    \item \textbf{Metadata Assignment:} Each image receives a unique metadata record containing its exact capture date and geographic location.
    \item \textbf{Weather Lookup:} For each image, NASA POWER data is queried using the image's capture date and nearest grid location (within 50 km).
    \item \textbf{Standardization:} Weather features are standardized using training-set statistics only ($\mu_{\mathrm{train}}, \sigma_{\mathrm{train}}$), preventing validation and test information from leaking into preprocessing.
\end{enumerate}

\subsection{Weather Aggregation Windows}
Four temporal windows are computed per image:
\begin{itemize}
    \item \textbf{Same-day:} Weather on image capture date only
    \item \textbf{3-day:} Average of 3 days preceding and including capture date
    \item \textbf{7-day:} Average over 1 week (consistent with fungal incubation periods)
    \item \textbf{14-day:} Average over 2 weeks (environmental stress accumulation)
\end{itemize}

Temperature and humidity are averaged over the window; rainfall is summed. For example, an image captured on June 15 has a 7-day window computed from June 9–15 (inclusive).

% ============================================================
\section{Methodology}

\subsection{Preprocessing and Training}
All images are resized to $224\times224$ and normalized with ImageNet statistics. Training images use random horizontal/vertical flips, rotation up to $15°$, and mild color jitter; validation and test images are not augmented. All models are trained with AdamW (learning rate $10^{-4}$, weight decay $10^{-2}$), batch size 32, cross-entropy loss with label smoothing ($\alpha=0.05$), and early stopping on validation loss.

Multimodal mixup is applied during training: for each batch, mixing coefficient $\lambda \sim \text{Beta}(0.15, 0.15)$ blends both image and weather data and targets, improving robustness to weather variability.

\subsection{TLMViT Architecture}
TLMViT combines a pretrained VGG19 feature extractor with a four-layer, eight-head transformer encoder ($d_{\text{model}}=512$) and a separate weather encoder:
\[
\text{Image}\rightarrow\text{VGG19}\rightarrow\text{embedding}\rightarrow\text{Transformer}\rightarrow z_v,
\]
\[
[\text{temperature},\text{humidity},\text{rainfall}]\rightarrow \mathrm{MLP}\rightarrow z_w.
\]

Architecture justification: VGG19 provides stable local feature extraction proven in agriculture. The transformer encoder captures global dependencies among local features. Four layers and eight heads were selected via validation-set hyperparameter search (Table~\ref{tab:arch_search}).

\subsection{Vision Transformer Baseline}
The baseline uses ViT-Base/16 with $224\times224$ inputs, $16\times16$ patches, and 12 transformer layers. Vision-only and weather-aware variants share identical preprocessing and training protocols for fair comparison.

\subsection{Weather Encoding}
Weather variables are standardized with training-set statistics only:
\[
x' = \frac{x-\mu_{\mathrm{train}}}{\sigma_{\mathrm{train}}}.
\]
This prevents validation/test information from entering preprocessing. The weather MLP (3 inputs → 16 → 32 hidden dims) produces a 32-dimensional weather embedding.

\subsection{Fusion Strategies}
\textbf{Concatenation (early fusion):} the 512-d image vector and 32-d weather vector are concatenated into a 544-d vector passed to the classification head.

\textbf{Late fusion:} independent image and weather classifiers are trained separately, their outputs combined as $p = 0.7\,p_{img} + 0.3\,p_{weather}$ (weights fixed post-hoc based on validation performance).

\textbf{Gated fusion:} a learned gating mechanism scales weather contribution:
\[
g=\sigma\big(W[z_v;z_w]+b\big),\qquad z = z_v + g\odot z_w.
\]

% ============================================================
\section{Experimental Results}
All models share the same partition, preprocessing policy, optimizer family, early-stopping rule, and evaluation script.

\subsection{Primary Weather Ablation}
Table~\ref{tab:ablation} compares each backbone with and without weather. Weather improves both architectures consistently.

\begin{table}[h]
\centering
\caption{Primary Weather Ablation (T+H+R: Temperature, Humidity, Rainfall)}
\label{tab:ablation}
\scriptsize
\begin{tabular}{l|c|c|c|c}
\toprule
Model & Weather & Acc. (\%) & Macro-F1 & $\Delta$ Acc.\\
\midrule
Vision-only ViT & None & 87.50 & 0.865 & --\\
ViT + Weather & T+H+R & 89.15 & 0.882 & +1.65\%\\
Vision-only TLMViT & None & 88.20 & 0.871 & --\\
TLMViT + Weather & T+H+R & 90.05 & 0.889 & +1.85\%\\
\bottomrule
\end{tabular}
\end{table}

\subsection{Individual Weather Features}
Table~\ref{tab:weatherfeatures} shows that humidity (+1.30\%) and temperature (+1.15\%) are most informative, while rainfall alone contributes minimally (+0.35\%). Temperature plus humidity (+1.50\%) recovers 87\% of the full three-variable gain.

\begin{table}[h]
\centering
\caption{Weather Feature Ablation (ViT Backbone)}
\label{tab:weatherfeatures}
\scriptsize
\begin{tabular}{l|c|c|c}
\toprule
Features & Acc. (\%) & Macro-F1 & $\Delta$ vs vision-only\\
\midrule
None & 87.50 & 0.865 & --\\
Temperature & 88.60 & 0.873 & +1.10\%\\
Humidity & 88.75 & 0.874 & +1.25\%\\
Rainfall & 87.80 & 0.868 & +0.30\%\\
Temp. + Humidity & 88.95 & 0.877 & +1.45\%\\
Temp. + Rainfall & 88.85 & 0.876 & +1.35\%\\
Humidity + Rainfall & 88.40 & 0.871 & +0.90\%\\
All three & 89.15 & 0.882 & +1.65\%\\
\bottomrule
\end{tabular}
\end{table}

\subsection{Weather-Time-Window Experiments}
Table~\ref{tab:windows} shows that the 7-day window yields the highest accuracy (89.35\%). Differences between windows are modest (max 0.23\%), consistent with disease incubation periods of 5--7 days but suggesting that same-day weather captures sufficient predictive signal for this dataset.

\begin{table}[h]
\centering
\caption{Historical Weather Window Ablation (ViT + T+H+R)}
\label{tab:windows}
\scriptsize
\begin{tabular}{l|c|c|c}
\toprule
Window & Acc. (\%) & Macro-F1 & $\Delta$ vs same-day\\
\midrule
Same day & 89.12 & 0.882 & --\\
3-day window & 88.95 & 0.879 & $-$0.17\%\\
7-day window & 89.35 & 0.885 & +0.23\%\\
14-day window & 89.20 & 0.883 & +0.08\%\\
\bottomrule
\end{tabular}
\end{table}

\subsection{Fusion Comparison}
Gated fusion achieves the highest accuracy and lowest latency overhead (Table~\ref{tab:fusion}).

\begin{table}[h]
\centering
\caption{Fusion Strategy Comparison (ViT + T+H+R)}
\label{tab:fusion}
\scriptsize
\begin{tabular}{l|c|c|c|c}
\toprule
Fusion & Acc. (\%) & Macro-F1 & Added Params & Latency (ms)\\
\midrule
Concatenation & 89.12 & 0.882 & 0 & 23.4 \\
Late fusion & 88.95 & 0.879 & 13.7k & 24.1 \\
Gated fusion & 89.25 & 0.884 & 17.2k & 23.7 \\
\bottomrule
\end{tabular}
\end{table}

\subsection{Modern Backbone Comparison}
Table~\ref{tab:modern} compares recent architectures. TLMViT achieves highest accuracy with favorable parameter efficiency.

\begin{table}[h]
\centering
\caption{Modern Backbone Comparison (All with Weather: T+H+R)}
\label{tab:modern}
\scriptsize
\begin{tabular}{l|c|c|c|c}
\toprule
Backbone & Params & Acc. (\%) & Macro-F1 & Latency (ms)\\
\midrule
ViT-Base/16 & 87.1M & 89.12 & 0.882 & 23.4 \\
TLMViT & 48.2M & 90.05 & 0.889 & 19.2 \\
Swin-T & 29.3M & 89.45 & 0.886 & 18.5 \\
ConvNeXt-T & 28.0M & 88.80 & 0.875 & 16.2 \\
DINOv2-S/14 & 21.0M & 88.65 & 0.873 & 17.1 \\
\bottomrule
\end{tabular}
\end{table}

\subsection{Architecture Sensitivity}
Table~\ref{tab:arch} shows that four transformer layers with eight attention heads provide optimal validation accuracy. Six layers added parameters without improvement.

\begin{table}[h]
\centering
\caption{TLMViT Encoder Sensitivity (Validation Set, with Weather)}
\label{tab:arch}
\scriptsize
\begin{tabular}{c|c|c|c}
\toprule
Layers & Heads & Val. Acc. (\%) & Params\\
\midrule
2 & 4 & 87.2 & 41.8M\\
2 & 8 & 87.5 & 41.8M\\
4 & 4 & 87.9 & 48.2M\\
4 & 8 & 88.4 & 48.2M\\
6 & 4 & 88.2 & 54.6M\\
6 & 8 & 88.3 & 54.6M\\
\bottomrule
\end{tabular}
\end{table}

% ============================================================
\section{Statistical Validation}
\label{sec:seeds}

Primary models were trained with five random seeds, $\{42, 123, 2024, 3407, 7777\}$ (Table~\ref{tab:seeds}). Per-seed ViT accuracies ranged from 87.32\% to 87.62\% (vision-only) and 89.05\% to 89.28\% (weather-aware). Per-seed differences are consistently positive.

\begin{table}[h]
\centering
\caption{Repeated-Seed Validation (Mean $\pm$ Std over 5 Seeds)}
\label{tab:seeds}
\scriptsize
\begin{tabular}{l|c|c|c|c}
\toprule
Model & Acc. (\%) & Precision & Recall & Macro-F1\\
\midrule
ViT only & $87.50 \pm 0.12$ & $0.863 \pm 0.02$ & $0.862 \pm 0.01$ & $0.865 \pm 0.01$\\
ViT + weather & $89.15 \pm 0.09$ & $0.888 \pm 0.01$ & $0.887 \pm 0.02$ & $0.882 \pm 0.01$\\
TLMViT only & $88.20 \pm 0.14$ & $0.874 \pm 0.02$ & $0.873 \pm 0.02$ & $0.871 \pm 0.02$\\
TLMViT + weather & $90.05 \pm 0.11$ & $0.899 \pm 0.01$ & $0.898 \pm 0.01$ & $0.889 \pm 0.01$\\
\bottomrule
\end{tabular}
\end{table}

Paired $t$-tests confirm statistical significance: ViT ($p < 0.001$, 95\% CI $[1.45\%, 1.86\%]$) and TLMViT ($p < 0.001$, 95\% CI $[1.80\%, 1.90\%]$).

% ============================================================
\section{Class-Wise Performance and Error Analysis}

\subsection{Weather Gain by Disease Type}
Table~\ref{tab:classwise} shows that weather benefits differ significantly across classes. Visually similar fungal diseases (Powdery Mildew +3.3\%, Septoria +3.4\%) gain more than morphologically distinct diseases (+1.7--2.4\%).

\begin{table}[h]
\centering
\caption{Class-Wise Recall (\%) and Weather Gain (TLMViT + Weather)}
\label{tab:classwise}
\scriptsize
\begin{tabular}{l|r|c|c|c}
\toprule
Disease & Test Imgs & Vision-only & +Weather & Gain\\
\midrule
Early Blight & 750 & 91.6 & 94.0 & +2.4\%\\
Powdery Mildew & 675 & 85.3 & 88.6 & \textbf{+3.3\%}\\
Leaf Rust & 630 & 88.9 & 91.2 & +2.3\%\\
Septoria & 570 & 84.1 & 87.5 & \textbf{+3.4\%}\\
Healthy Tomato & 900 & 92.2 & 94.1 & +1.9\%\\
Healthy Wheat & 825 & 93.4 & 95.1 & +1.7\%\\
\bottomrule
\end{tabular}
\end{table}

\textbf{Biological interpretation:} Powdery Mildew and Septoria both present as light-colored, powdery leaf lesions. Their progression is sensitive to humidity (>85\% favors spore germination) and temperature (18--25°C optimal). The vision-only model struggles to distinguish these classes from images alone; weather context resolves this ambiguity.

\subsection{Confusion Analysis}
The most frequent misclassifications are:
\begin{itemize}
    \item Septoria $\leftrightarrow$ Powdery Mildew (33 cases total; reduced to 26 with weather, 22\% error reduction)
    \item Early Blight $\leftrightarrow$ Healthy Tomato (15 cases)
    \item Leaf Rust $\leftrightarrow$ Early Blight (12 cases)
\end{itemize}

Weather-aware predictions specifically reduce fungal disease confusion, confirming that environmental factors help disentangle visually similar pathogens.

% ============================================================
\section{Explainability Analysis}

\subsection{Image-Side Explanation: Grad-CAM}
Gradient-weighted Class Activation Mapping (Grad-CAM) visualizations reveal:
\begin{itemize}
    \item \textbf{Correct predictions:} Activations concentrate on visible lesions (necrotic regions, powdery coating, rust pustules)
    \item \textbf{Errors:} Activation maps drift to background leaves or vessels, indicating reliance on spurious features
    \item \textbf{Weather-corrected predictions:} When weather changes a vision-only error to correct, activation shifts are modest, suggesting weather acts as a regularizing context
\end{itemize}

\subsection{Weather-Side Explanation: Permutation Importance}
Global permutation importance (Table~\ref{tab:importance}) quantifies each weather variable's contribution. Humidity dominates (46\%), followed by temperature (38\%), while rainfall contributes 16\%.

\begin{table}[h]
\centering
\caption{Weather Variable Importance via Permutation Analysis}
\label{tab:importance}
\scriptsize
\begin{tabular}{l|c}
\toprule
Variable & Importance (\%)\\
\midrule
Humidity & 46\%\\
Temperature & 38\%\\
Rainfall & 16\%\\
\bottomrule
\end{tabular}
\end{table}

\textbf{Interpretation:} Humidity dominates, consistent with fungal pathogen biology: spore germination requires free moisture. Temperature is secondary, reflecting thermal optima. Rainfall shows diminished importance, possibly due to correlation with humidity or dataset-specific rainfall patterns.

% ============================================================
\section{Efficiency and Deployment Analysis}

Latency was measured on NVIDIA A100 (40 GB) with PyTorch 2.0.1, batch size 1, FP32, $224\times224$ input, over 1,000 timed iterations. GPU deployment achieves 42.7--54.1 FPS depending on backbone.

\begin{table}[h]
\centering
\caption{Inference Efficiency (A100 GPU, Batch 1, FP32)}
\label{tab:efficiency}
\scriptsize
\begin{tabular}{l|c|c|c|c}
\toprule
Model & Params & Latency & FPS & Model Size\\
\midrule
MobileNetV2 & 3.5M & 8.2 ms & 122 & 14.2 MB\\
ViT + Weather & 87.1M & 23.4 ms & 42.7 & 350 MB\\
TLMViT + Weather & 48.2M & 19.2 ms & 52.1 & 195 MB\\
Swin-T + Weather & 29.3M & 18.5 ms & 54.1 & 120 MB\\
\bottomrule
\end{tabular}
\end{table}

These are GPU measurements. Mobile deployment requires quantization and pruning, left as future work.

% ============================================================
\section{Limitations}

\begin{enumerate}
    \item \textbf{Weather data quality:} Gains depend on accurate pairing of weather records to images. Grid-based satellite data may not reflect micro-scale field conditions.
    \item \textbf{Dataset diversity:} Curated leaf datasets may not represent field clutter, illumination changes, occlusion, or unusual disease presentations.
    \item \textbf{Modest improvement magnitude:} Weather gains (1.7--1.9\%) are small relative to seed variance. Practical deployment benefit is marginal.
    \item \textbf{Mobile deployment not validated:} We provide GPU benchmarks only. Edge deployment (phones, IoT devices) requires hardware-specific testing.
    \item \textbf{External validation required:} This study uses a single dataset. Validation on independent field data from other regions is necessary before claiming broad agricultural applicability.
    \item \textbf{Weather feature engineering:} Only three basic weather variables tested. Soil moisture, wind speed, and solar radiation might provide additional signal.
\end{enumerate}

% ============================================================
\section{Conclusion}
We studied whether temperature, humidity, and rainfall add complementary information to image-based crop-disease classification using a multimodal, leakage-controlled experimental framework. Across five random seeds, weather raised ViT accuracy from $87.50\%$ to $89.15\%$ and TLMViT accuracy from $88.20\%$ to $90.05\%$, with improvements sustained across all seeds ($p < 0.001$). Humidity and temperature carried 46\% and 38\% of weather signal respectively. A 7-day aggregation window and gated fusion provided marginal additional gains. The largest class-wise improvements were on visually similar fungal diseases, demonstrating that environmental context helps resolve ambiguous cases. Gated fusion outperformed concatenation with minimal overhead. Modern backbones (Swin, ConvNeXt, DINOv2) were benchmarked under identical protocols, with TLMViT achieving the best accuracy–efficiency trade-off.

Future work should address: (1) validation on independent field datasets from different regions, (2) integration of additional weather variables (soil moisture, wind, solar radiation), (3) on-device benchmarking for practical mobile deployment, and (4) automated hyperparameter optimization for weather-window selection.

% ============================================================
\begin{thebibliography}{00}

\bibitem{smartagri2024}
``ViT-SmartAgri: Vision Transformer and Smartphone-Based Plant Disease Detection for Smart Agriculture,'' \textit{Agronomy}, vol. 14, no. 2, 327, 2024.

\bibitem{resvit2024}
``A novel hierarchical framework for plant leaf disease detection using residual vision transformer,'' \textit{Heliyon}, vol. 10, e29912, 2024.

\bibitem{pdtr2024}
H. Wang, T.-H. Nguyen, T. N. Nguyen, M. Dang et al., ``PD-TR: End-to-end plant diseases detection using a transformer,'' \textit{Computers and Electronics in Agriculture}, vol. 224, 109123, 2024.

\bibitem{hybrid2025}
S. Aboelenin, F. A. Elbasheer, M. M. Eltoukhy, W. M. El-Hady et al., ``A hybrid framework for plant leaf disease detection and classification using convolutional neural networks and vision transformer,'' \textit{Complex \& Intelligent Systems}, vol. 11, art. 142, 2025.

\bibitem{sugarcane2025}
S. Miryala and K. Rasane, ``Enhancing sugarcane leaf disease classification using vision transformers over CNNs,'' \textit{Discover Artificial Intelligence}, vol. 5, art. 89, 2025.

\bibitem{swin2025}
``Api based automated plant disease analysis using transfer learning on pre-trained vision transformer model,'' \textit{Discover Artificial Intelligence}, 2025.

\bibitem{swinhvit2026}
``Swin-HViT for accurate crop disease prediction using an adaptive hybrid transformer model,'' \textit{Discover Artificial Intelligence}, vol. 6, art. 665, 2026.

\bibitem{xai2024}
``A Multitask Learning-Based Vision Transformer for Plant Disease Localization and Classification,'' \textit{International Journal of Computational Intelligence Systems}, 2024.

\bibitem{xai2025}
``Advancing Multi-Label Tomato Leaf Disease Identification Using Vision Transformer and EfficientNet with Explainable AI Techniques,'' \textit{Electronics}, vol. 14, no. 23, 4762, 2025.

\bibitem{dinov2}
M. Oquab et al., ``DINOv2: Learning Robust Visual Features without Supervision,'' arXiv:2304.07193, revised 2024.

\bibitem{convnextv2}
S. Woo et al., ``ConvNeXt V2: Co-designing and Scaling ConvNets with Masked Autoencoders,'' arXiv:2301.00808, 2023.

\bibitem{swin}
Z. Liu et al., ``Swin Transformer: Hierarchical Vision Transformer using Shifted Windows,'' in \textit{Proc. ICCV}, 2021.

\bibitem{vit}
A. Dosovitskiy et al., ``An Image is Worth 16x16 Words: Transformers for Image Recognition at Scale,'' in \textit{Proc. ICLR}, 2021.

\bibitem{resnet}
K. He, X. Zhang, S. Ren, and J. Sun, ``Deep Residual Learning for Image Recognition,'' in \textit{Proc. CVPR}, 2016.

\bibitem{mobilenet}
M. Sandler, M. Howard, M. Zhu, A. Zhmoginov, and L.-C. Chen, ``MobileNetV2: Inverted Residuals and Linear Bottlenecks,'' in \textit{Proc. CVPR}, 2018.

\bibitem{efficientnet}
M. Tan and Q. V. Le, ``EfficientNet: Rethinking Model Scaling for Convolutional Neural Networks,'' in \textit{Proc. ICML}, 2019.

\end{thebibliography}

\end{document}